\DocumentMetadata{lang=en-US, pdfstandard=ua-2, testphase={phase-III,table}}
\documentclass{./template/articlex}

\hypersetup{pdftitle={Why Economics? What UL Lafayette Graduates Actually Earn},
            pdflang={en-US}}

% --- one-pager formatting -----------------------------------------------
\usepackage{geometry}
\geometry{landscape, left=0.55in, right=0.55in, top=0.5in, bottom=0.45in}
\usepackage{xcolor}
\usepackage{tcolorbox}
\usepackage[scaled=0.92]{helvet}

\definecolor{ULRED}{HTML}{CE0E2D}
\definecolor{INK}{HTML}{222222}
\definecolor{SOFT}{HTML}{555555}

% hero number with caption underneath
\newcommand{\stat}[2]{%
  \begin{minipage}[t]{0.315\linewidth}
    \vspace{0pt}
    {\sffamily\bfseries\color{ULRED}\fontsize{24}{26}\selectfont #1}\par\vspace{1pt}
    {\sffamily\color{SOFT}\fontsize{8.6}{10.2}\selectfont #2}
  \end{minipage}}

% a job title in the careers panel
\newcommand{\job}[1]{{\sffamily\color{INK}\fontsize{9}{12}\selectfont #1}\par}

\setlength{\parskip}{0pt}
\setlength{\parindent}{0pt}
\pagestyle{empty}
% ------------------------------------------------------------------------

\begin{document}
\begin{singlespace}

%% ---------- header -----------------------------------------------------
{\sffamily\bfseries\color{ULRED}\fontsize{23}{26}\selectfont
Why Economics? What UL Lafayette Graduates Actually Earn}

\vspace{2pt}
{\color{ULRED}\rule{\linewidth}{1.5pt}}
\vspace{4pt}

{\sffamily\color{SOFT}\fontsize{10.5}{13}\selectfont
The real paychecks of real UL graduates. The U.S.\ Census Bureau matches UL
Lafayette's own transcript records to the earnings records employers file for
their workers --- so these are not survey answers, and nobody is missing because
they declined to reply.}

\vspace{8pt}

%% ---------- three numbers ----------------------------------------------
\noindent
\stat{\$105,369}{Median earnings of UL Economics graduates ten years after they
finish their degree, in January 2026 dollars.}
\hfill
\stat{Highest}{Of the seven undergraduate business majors at UL Lafayette, at
ten years out --- at the median and at the 75th percentile alike.}
\hfill
\stat{+48\%}{More than the typical UL Lafayette graduate across every field on
campus, at the ten-year median.}

\vspace{10pt}

%% ---------- figure beside careers ---------------------------------------
\noindent
\begin{minipage}[t]{0.665\linewidth}
  \vspace{0pt}
  \includegraphics[width=\linewidth, alt={Bar chart of median and
  75th-percentile earnings for seven UL Lafayette undergraduate business majors
  at five and ten years after graduation. Economics leads both measures at ten
  years, at \$105,369 median and \$141,919 at the 75th percentile; Accounting is
  next at \$85,829 median.}]%
  {figures/econ_vs_business_majors_5yrs_10yrs_2026-09-28.pdf}

  \vspace{4pt}

  {\sffamily\color{SOFT}\fontsize{8.3}{10.4}\selectfont
  \textbf{\color{INK}Median} is the graduate in the middle --- half earned more,
  half earned less. \textbf{\color{INK}75th percentile} is the graduate
  three-quarters of the way up: one in four earned at least that much. At five
  years the business majors look much alike; the gap opens later, which is why
  the ten-year numbers are the ones worth reading.}
\end{minipage}
\hfill
\begin{minipage}[t]{0.305\linewidth}
  \vspace{0pt}
  \begin{tcolorbox}[colback=ULRED!4, colframe=ULRED, boxrule=0.9pt, arc=2pt,
                    halign=flush left,
                    left=8pt, right=8pt, top=6pt, bottom=6pt]
  {\sffamily\bfseries\color{ULRED}\fontsize{11}{13}\selectfont
  Where Economics majors actually work}

  \vspace{3pt}

  {\sffamily\color{SOFT}\fontsize{8.6}{10.6}\selectfont\raggedright
  An Accounting major applies for jobs called ``Accountant.'' Economics does not
  work that way, and that is the biggest reason students underestimate it.
  Almost none of these titles contain the word \emph{economist}. All of them
  hire economics majors.\par}

  \vspace{6pt}

  \job{Financial Analyst}
  \job{Credit Analyst / Underwriter}
  \job{Pricing \& Revenue Analyst}
  \job{Data \& Business Analyst}
  \job{Market Research Analyst}
  \job{Management Consultant}
  \job{Policy Analyst}
  \job{Economic Consultant}
  \job{Research Assistant --- Fed, Census, BLS}

  \vspace{5pt}

  {\sffamily\color{SOFT}\fontsize{8.4}{10.4}\selectfont\raggedright
  Employers are buying the ability to work out what is driving a number and what
  would happen if something changed --- then explain it to someone who does not
  want the math. Economics is also a well-worn path into law school.\par}
  \end{tcolorbox}
\end{minipage}

\vspace{7pt}

%% ---------- footer ------------------------------------------------------
{\sffamily\color{SOFT}\fontsize{7.2}{8.6}\selectfont\raggedright
\textbf{Economics is the red bar}; majors are ordered by ten-year median
earnings. Covers bachelor's graduates, classes of 2001--2021, who were working
steadily in the year measured --- Census excludes anyone earning less than
full-time minimum wage or with two or more quarters of no earnings. Census
reports program codes finer than UL's degrees, so related codes are grouped
(Management Information Systems within Management, Taxation within Accounting)
and weighted by graduate counts; Economics is a single code, so its bars are
exactly what Census published. Graduates who later earned an MBA, law or
master's degree still count here, and double majors appear in both bars. These
figures show what graduates of each major earned, not what the major did for
them. \textbf{Source:} U.S.\ Census Bureau, Longitudinal Employer--Household
Dynamics Program, Post-Secondary Employment Outcomes (PSEO), Louisiana file,
\url{https://lehd.ces.census.gov/data/pseo_experimental.html}; accessed
September 28, 2026. January 2026 dollars, adjusted with the CPI-U (BLS series
CUUR0000SA0).\par}

\end{singlespace}
\end{document}